\chapter{Current limitations}
\label{chap:limitations}

\section{Thin vessels are below the resolution}
\label{sec:thin}

With a blur of $\sigma = 4.4$\um{}, the image of a point is about 10\um{} wide at half its height. A
vessel of 3 to 6\um{} radius is narrower than the blur: its image is essentially the blur's own shape,
dimmer than a large vessel, whatever its true size. On the phantom, vessels of that size are found
\pbValueVesselRecallThin{} of the time, against \pbValueVesselRecallMedium{} for 6--12\um{} and
\pbValueVesselRecallLarge{} above.

I tried two recoveries and dropped them. The missed thin vessels turned out not to be faint at all (their
voxels are mostly brighter than the bone threshold): they are lost because, next to a brighter vessel,
they fall below that vessel's half-maximum and are taken for its halo. Telling them apart from a real
halo with the Hessian recovered about as many wrong voxels as right ones: no net gain, so I did not
keep the code. The
consequence for future measurements is that diameters under about 12\um{} will be overestimated, and
the thinnest vessels partly missed: this is a property of the scan, not of the software.

\section{Rods and vessels}

Shape alone cannot tell a trabecular rod from a vessel: both are tubes. The decision rests on what they
touch in 3D. A dim vessel running along a trabecula, or a rod touching vessels on all sides, can still
be misclassified; these are the same ``zones of uncertainty'' we met in 2018.

\section{Validation}

The phantom is calibrated on the real scan, but it is synthetic: its vessel paths, burst textures and
bone geometry are simplifications. The scores of Chapter~\ref{chap:validation} measure the method on
realistic data, not the segmentation of this particular knee. Expert labels on a few real crops (the
annotation page is ready) would close that gap.

\section{The data}

\begin{itemize}
  \item \textbf{One sample.} I could only check the parameters on a single scan. Physical units make them
        transferable in principle, not in fact.
  \item \textbf{Voxel size.} 3\um{} is the slice thickness we used in 2018 and agrees with the
        anatomy; I assumed isotropy.
  \item \textbf{A thin slab.} The stack is 100 slices thick (0.3\,mm) for a knee several millimetres
        wide: most vessels leave the volume, which limits what 3D connectivity can decide near the
        first and last slices.
  \item \textbf{Bursts and vessels.} A vessel running into a burst cluster is partly absorbed by it.
\end{itemize}

\section{The original goal}

The brief asked for a quantified network: nodes, branches, calibre, branch lengths. pyBones delivers a
validated segmentation and meshes, the prerequisite, but not the quantification yet.
